\chapter{复利的真相，和关于复利的谎言}
\chaptermeta{22}

\begin{ledetext}
复利被讲成了理财界的鸡汤。它是真的，但它有五个前提，缺一个就不成立。而它对亏损的态度，比你想的狠得多。
\end{ledetext}

\section{先做一道小学算术}

你有 10 万块。第一年亏 50\%，第二年赚 50\%。

现在你有多少？

7 万 5。

10 万亏一半剩 5 万，5 万涨一半是 7 万 5。两年下来"平均收益率"是 0\%，你的钱少了四分之一。要回到 10 万，第二年得赚 100\%，不是 50\%。

你的账户不做加法，它只做乘法。所谓复利，算的是\term{几何平均}{geometric mean}\index{026@几何平均}。而几何平均对亏损的惩罚是不对称的：\textbf{亏得越深，回本所需的涨幅涨得越快，而且是加速涨。}

这张表建议你拍下来。

\begin{tablewrap}
\begin{xltabular}{\linewidth}{>{\raggedright\arraybackslash}p{0.20\linewidth}Xrr}
\toprule
\thead{你亏了} & \thead{剩下} & \thead{回本需要涨} & \thead{按年化 8\% 算，要爬多少年} \\
\midrule
\textbf{10\%} & 0.90 & 11.1\% & 1.4 年 \\
\textbf{20\%} & 0.80 & 25.0\% & 2.9 年 \\
\textbf{30\%} & 0.70 & 42.9\% & 4.6 年 \\
\textbf{50\%} & 0.50 & 100\% & 9.0 年 \\
\textbf{70\%} & 0.30 & 233\% & 15.6 年 \\
\textbf{90\%} & 0.10 & 900\% & 29.9 年 \\
\bottomrule
\end{xltabular}
\end{tablewrap}

最后一行读三遍。亏掉九成之后，就算你接下来每年都能稳拿 8\%，回到原点也要将近三十年。三十年是很多人的整个工作生涯。

\begin{egbox}{举个栗子}
汪工在成都一家做汽车线束的厂子里当工艺工程师。2015 年 6 月，他把攒了六年的 37 万 8 全部买进一只创业板主题基金，因为"身边人都在赚"。

后面的事你知道。他没割，一直拿着，最深的时候账户剩 14 万 3，亏掉 62\%。

62\% 意味着什么？意味着他需要涨 163\% 才回到 37 万 8。按每年 8\% 爬，要爬 12 年 7 个月。他那年 34 岁。

这不是"拿得住就行"能解决的问题。这是算术。

\end{egbox}

\section{8\% 和 8\%，不是同一个 8\%}

两个人都告诉你，他的长期年化是 8\%。

第一个是真的每年 8\%，稳得像一张持有到期的国债。第二个是隔年跳：今年 +30\%，明年 −14\%，来回摆。你把 30 和 −14 一加除以二，正好是 8。

各拿 32 万，各走三十年。

\begin{barfig}
  \barrow{每年稳稳 8\%}{100}{322 万}
  \barrow{+30\% 与 −14\% 交替}{53}{170 万}
\end{barfig}

差 151 万。而这两个人的"平均年化"在任何一份宣传材料上都会被写成同一个数字。

波动的那位，真实年化是 5.74\%，不是 8\%。因为 1.30 乘 0.86 等于 1.118，开平方是 1.0574。少掉的那 2.26 个百分点，被波动本身吃掉了。这个东西叫\term{波动拖累}{variance drag}\index{008@波动拖累}，有个粗略但好用的口算公式：\textbf{几何收益率约等于算术收益率减去方差的一半。}上面这个例子波动率是 22 个百分点，0.22 的平方除以 2 是 0.0242，8\% 减 2.42\% 得 5.58\%，跟真实值 5.74\% 只差一点点。

把它想成开车：同样从成都到重庆，一个走高速匀速跑，一个在山路上不停地急刹再地板油。里程一样，油耗差一截。可以继续推下去：既然波动本身就是成本，那么在收益率不变的前提下把波动压下来，等于凭空提高了收益。这正是第 21 章说资产配置是"免费午餐"的真正原因，它不是玄学，是这条公式。

\section{钱几乎全在后半段}

有个心算工具叫 72 法则：用 72 除以年化收益率，得到本金翻一倍需要的年数。6\% 就是 12 年，9\% 就是 8 年。误差在常见区间内很小。

拿 23 万，按 6\% 走。第 12 年 46 万 3，第 24 年 93 万 1，第 30 年 132 万 1。

现在把这条曲线切成两段看。

\begin{statrow}{2}
  \statitem{50.8}{万}{前 20 年一共长出来的钱}
  \statitem{58.3}{万}{第 21 到第 30 年长出来的钱}
\end{statrow}

最后十年的收成，超过了前二十年的总和。

复利的形状不是一条斜坡，是一根被慢慢压弯的鱼竿。前二十年你几乎看不出它在弯，最后十年它突然抬起来。这个形状能推出两个很实际的结论。

第一，早开始比高收益更值钱。

\begin{egbox}{举个栗子}
小林 26 岁，郑州一家社区医院的护士，每月定投 1800 块，一年 21600。她只投十年，36 岁那年停手，之后一分钱不再加，就那么放着。按 6\% 算，她 61 岁时账上大概 122 万。

她表哥 36 岁才开始，同样每月 1800，一直投到 61 岁，投了 25 年，本金掏出去 54 万。他 61 岁时账上大概 118 万 5。

小林掏了 21 万 6，表哥掏了 54 万，多掏了 32 万多。最后表哥还少几万。

差别不在收益率，两个人用的是同一个 6\%。差别在小林那笔钱多在曲线上待了十年，而那十年正好是鱼竿抬头的那一段。

\end{egbox}

第二，中途取钱的代价大得不像话。

还是那 23 万、那 6\%。假设第 12 年你抽走 15 万去付一笔首付。到第 30 年，账户不是少 15 万，是少 42 万 8。因为被抽走的不只是那 15 万，还有它后面 18 年的弯曲。

那 15 万的真实价格，是 42 万 8。

\begin{pullquote}{}
复利不奖励聪明。它奖励活得久、开始得早、以及最不像本事的那一项：不动手。
\end{pullquote}

\section{榜单上看不见的那些基金}

第一个谎言：把"历史年化 X\%"当成承诺。

"这只基金过去五年年化 19\%。"这句话通常是真的，它同时几乎是无用的。

排行榜是从\emph{还活着}的产品里排出来的。清盘的、被合并的、换了基金经理又换了名字重新出发的，都不在榜上。中国公募基金规模 2026 年初约 38 万亿元，几千只产品，每年都有一批悄悄退场。你看到的那张榜单，是一张幸存者的合影，照片里没有没能来的人。这叫\term{幸存者偏差}{survivorship bias}\index{073@幸存者偏差}。

更大的那个幸存者，是国家。

"美股长期年化 10\%"这句话背后，站着一个本土从未被现代战火犁过、打赢了两次世界大战、货币成了全世界储备资产的国家。同一段时间里，1917 年的俄国股市、1949 年的上海证券交易所、两次大战之间中欧的一批市场，投资者最终拿到的数字是零。

\textbf{把一个赢家的历史收益率，当成"资产类别的自然属性"，这个推理是错的。}它不是"股票长期会涨 10\%"，它是"在过去一百年里赢麻了的那个国家，股票涨了 10\%"。这两句话的区别，就是你未来三十年计划里最大的那个假设。

\begin{deepbox}{进阶：可以跳过}
归零的市场有个讨厌的性质：它们不进入任何长期收益率序列。因为序列需要数据，而归零之后就没有数据了。所以所有"长期看股票一定涨"的统计，天然只能用活下来的样本做。

这不是说统计没用。是说你在读那个数字的时候，得知道分母被谁悄悄改小了。

\end{deepbox}

顺便说说第三个谎言，因为它是同一个毛病：\textbf{"巴菲特靠的是复利。"}

他确实靠复利。但有两样东西很少被一起提。

一是杠杆。伯克希尔的核心是保险，保户交了保费、还没出险的那段时间，钱归公司支配，这笔钱叫\term{浮存金}{float}\index{016@浮存金}。它的成本常年低于国债利率，某些年份甚至是负的。学术界对他长期杠杆倍数的估算大约是 1.6 倍。用别人的钱、且几乎不要利息、且不会被追缴，这个条件本身就值一大截超额收益。

二是时间。他 11 岁买第一只股票，绝大部分财富是 50 岁之后堆起来的。他的复利跑道长度，是普通人的两倍还多。

说他"只靠复利"，是把一个幸存者的全部条件砍掉，只留下最容易复制的那一条给你看。

\section{复利的五个前提}

第二个谎言：复利适用于一切。

它需要同时满足五个条件。缺一个，公式就不成立了。

\begin{flowlist}
  \flowitem{01}{时间足够长}{至少二十年，最好三十年。如果你 47 岁才开始规划退休，鱼竿抬头的那一段你多半赶不上。这不是打击，是让你把预期挪到"提高储蓄率"上，那个变量还在你手里。}
  \flowitem{02}{中途不取钱}{结婚、买房、生孩子、父母住院、失业半年。普通人的三十年里，几乎必然有一次要动这笔钱。前面算过，第 12 年抽 15 万，终值少 42 万 8。}
  \flowitem{03}{收益能再投资}{分红要再投，票息要再投。拿去消费的那部分，不参与复利。一只年年高分红但你年年花掉的产品，你拿到的是单利。}
  \flowitem{04}{不破产}{任何一年出现 −100\%，前面所有年份的乘积一次性归零。加杠杆、押单一标的、玩合约的人，把这个条件的概率从"几乎不可能"提到了"迟早会碰上"。}
  \flowitem{05}{税和费吃不掉太多}{管理费、托管费、销售服务费、申赎费、换手带来的摩擦，加起来每年 1.5\% 不算夸张。8\% 变成 6.5\%，三十年后终值只剩原来的三分之二。你没做错任何事，钱少了三分之一。}
\end{flowlist}

\begin{mythbox}{常见误解}
\mvfrow{误解}{"我要求不高，年化 15\% 就行。"}
\mvfrow{事实}{2026 年 8 月，中国 7 天期逆回购政策利率 1.4\%，5 年期以上 LPR 3.5\%；美国联邦基金利率目标区间 3.50\%–3.75\%。要求 15\%，等于要求在无风险利率之上再拿 11 到 13 个百分点的补偿。这个补偿没有人会白给你，它对应的是你必须承担的风险。真正该问的不是"15\% 高不高"，而是"为了这 15\%，我要接受多大概率的一次 −40\%"。回到本章第一张表：−40\% 需要涨 66.7\% 才回本。}
\end{mythbox}

\section{同样的平均收益，倒过来就要命}

这一节讲一个中文世界里几乎没被讲清楚的东西：\term{序列风险}{sequence of returns risk}\index{074@序列风险}。

先说结论。在\textbf{积累期}，收益的先后顺序不重要，甚至前面跌还是好事，你每月定投的钱买到了更便宜的份额。在\textbf{提取期}，顺序决定生死。

郑老师，杭州人，教了三十二年中学化学，2026 年退休。账上 246 万，打算每年取 12 万补贴生活，取二十年。

给她两种未来。收益率的那二十个数字完全一样，只是顺序颠倒。

路径 A：头三年 −27\%、−17\%、−12\%，之后十七年每年 +10\%。\\ 路径 B：先是十七年每年 +10\%，最后三年 −27\%、−17\%、−12\%。

两条路的算术平均都是 5.7\%，几何年化都是 5.08\%，每年都取 12 万，二十年都取走 240 万。数学上，它们是同一组收益率。

\begin{tablewrap}
\begin{xltabular}{\linewidth}{>{\raggedright\arraybackslash}p{0.20\linewidth}rr}
\toprule
\thead{时点} & \thead{路径 A（先跌后涨）} & \thead{路径 B（先涨后跌）} \\
\midrule
\textbf{第 3 年末} & 105 万 & 284 万 \\
\textbf{第 10 年末} & 80 万 & 428 万 \\
\textbf{第 17 年末} & 31 万 & 708 万 \\
\textbf{第 19 年末} & 10 万 & 412 万 \\
\textbf{第 20 年末} & 0 & 352 万 \\
\bottomrule
\end{xltabular}
\end{tablewrap}

路径 B 走完二十年还剩 352 万，够她再活二十九年。路径 A 在第 19 年末只剩 10 万出头，第 20 年那笔 12 万她取不全。同样的平均收益，一个人留下三百多万，另一个人钱没了。

原因不复杂。头三年市场跌掉四成多，她还是得取 12 万，只能在低位卖出份额。卖掉的份额从此不在池子里，后面十七年的 +10\% 是给剩下的份额涨的。

像在旱季从水池里舀水。后来雨季来了，水位涨回去了，但你舀走的那几瓢不在池子里，它们没跟着涨。继续推一步：那么退休前后的三到五年，是一个人整个财务生命里风险最集中的窗口，因为它同时具备"本金最大"和"开始往外抽水"两个条件。所谓"临近退休降低权益仓位"，不是保守，是算术。

\begin{warnbox}{小心}
这一节不是教你择时。没有人知道你退休那年市场会涨还是会跌，任何声称知道的人都可以直接拉黑。

能做的是结构上的安排：让提取的那部分钱不依赖当年的市场。比如留出两到三年生活费的现金和短久期债券，跌的年份就动这一块，不动权益。这只是机制上的说明，不构成投资建议，具体怎么配得看你自己的现金流和承受力。

\end{warnbox}

\section{你只有一条时间线}

最后一件事，也是这本书里我最想让你带走的一个念头。

有个游戏。你押上全部身家抛一次硬币：正面，涨 50\%；反面，跌 40\%。

算期望收益：0.5 乘 1.5 加 0.5 乘 0.6，等于 1.05。每一轮期望赚 5\%。任何一本教科书都会告诉你这个游戏值得玩。

现在分两种玩法。

第一种：一百个人，每人玩一轮。这群人加总起来确实赚了大约 5\%，有人翻倍有人腰斩，平均是正的。

第二种：一个人，连玩一百轮。假设正反各出现五十次，他的 1 万块会变成 51 块 5。

因为 1.5 乘 0.6 等于 0.9。每凑齐一正一反，他就少掉 10\%。

\begin{statrow}{2}
  \statitem{131.5}{倍}{集合平均给出的答案（1.05 的 100 次方）}
  \statitem{0.5}{\%}{时间平均给出的答案（一个人玩 100 轮的剩余）}
\end{statrow}

这两个数字都没算错。它们回答的是两个不同的问题。前者叫\term{集合平均}{ensemble average}\index{025@集合平均}，后者叫\term{时间平均}{time average}\index{050@时间平均}。当两者不相等，这个系统就是\term{非遍历的}{non-ergodic}\index{014@非遍历的}。

\textbf{期望收益是替"很多个你"算的。而你只有一个你。}

\begin{pullquote}{}
一百个人各玩一次俄罗斯轮盘，这个游戏的期望收益是正的。一个人连玩一百次，他必死。
\end{pullquote}

所以在"把收益率从 8\% 提到 10\%"和"把爆仓概率从 3\% 降到 0"之间，永远先做后者。前者是把乘数从 1.08 改成 1.10，后者是保证这串乘数里不会出现一个 0。

任何数乘以 0 等于 0。小学二年级的内容，也是这本书里最贵的一句话。

\begin{mythbox}{常见误解}
\mvfrow{误解}{"拿得住就行，时间能抹平一切亏损。"}
\mvfrow{事实}{时间能抹平波动，抹不平永久性损失。退市的公司、违约清算的债、被强平的杠杆账户、清盘的基金，没有"后来"。就算不归零，时间也可能长到超出你的生命尺度：日经 225 在 1989 年 12 月 29 日收于 38915 点，再次站上这个数字是 2024 年，中间隔了三十四年；2026 年 6 月它第一次突破 70000 点。三十四年，是一个人从入职到退休的全程。}
\end{mythbox}

\begin{databox}{2026 数据：先把尺子摆好}
美国联邦基金利率目标区间 3.50\%–3.75\%，2026 年已连续五次维持不变（截至 7 月）。中国 7 天期逆回购政策利率 1.4\%；2026 年 8 月 20 日 LPR，1 年期 3.0\%，5 年期以上 3.5\%。

拿这几个数当基准线：在人民币资产里，任何明显高于 3.5\% 的"稳定收益"，高出来的部分都是风险的价格，不是谁的善意。这个换算你得会自己做，因为下一章里所有人都指望你不会做。

\end{databox}

\begin{takeaway}{本章一句话}
复利是乘法，不是加法。乘法怕两件事：怕中间有个特别小的数，更怕中间有个 0。

\begin{itemize}
  \item 亏 50\% 要涨 100\% 才回本；亏 90\% 按 8\% 爬要三十年。
  \item 波动本身在吃收益，压低波动等于提高收益。
  \item 复利的钱几乎全在最后十年，所以早开始比高收益重要，中途取钱代价极大。
  \item 提取期最怕开头就跌。同样的平均收益，顺序颠倒可以是"剩 352 万"和"归零"的差别。
  \item 先保证不出局，再谈收益率。你只有一条时间线。
\end{itemize}

\end{takeaway}

\begin{bridgebox}
\textbf{下一章}这一章讲的每一件事，骗子都比你更早学会了。他们知道你算不清几何平均，知道你分不清集合平均和时间平均，也知道你看到"稳定 15\%"时脑子里没有那把 3.5\% 的尺子。好消息是，所有骗局的零件加起来只有五个。
\end{bridgebox}

